\ChapterImageStar[cap:conclusiones]{Conclusiones}{./images/fondo.png}\label{cap:conclusiones}
\mbox{}\\

El desarrollo de esta investigación permitió alcanzar el objetivo general de definir aspectos de seguridad para la infraestructura de 
computación en la nube del grupo de investigación \GRID\ de la Universidad del Quindío. A través de un enfoque metodológico que integró la 
caracterización del entorno, el desarrollo de un Mapeo Sistemático de Literatura (\SMS), el análisis \DAR\ y el modelado arquitectónico 
mediante TOGAF y ArchiMate, fue posible identificar, analizar y seleccionar referentes metodológicos y tecnológicos alineados con las 
necesidades actuales de seguridad del entorno institucional.

El proceso de investigación permitió consolidar una arquitectura de defensa en profundidad orientada al monitoreo, correlación de eventos, 
detección de anomalías y respuesta ante incidentes de seguridad dentro de la infraestructura \GRID. Asimismo, la integración de tecnologías 
como Suricata, Wazuh y el SonicWall NSA 6600 permitió establecer un esquema interoperable de supervisión continua, fortaleciendo la 
trazabilidad operacional y la visibilidad sobre las actividades ejecutadas dentro de la infraestructura.

De igual forma, el trabajo desarrollado aporta una base reproducible y escalable para futuros procesos de implementación y fortalecimiento 
de arquitecturas de seguridad en entornos académicos y de investigación, proporcionando un modelo alineado con controles de la \ISO/IEC 
27001 e \ISO/IEC 27017 y respaldado por evidencia científica obtenida a partir del estudio sistemático realizado.

\section{Sobre la metodología empleada}
\noindent

La metodología desarrollada, compuesta por la caracterización del entorno \GRID, el desarrollo del Mapeo Sistemático de Literatura (\SMS) y 
la aplicación del análisis \DAR\ del modelo \CMMI, demostró ser un enfoque adecuado para la toma de decisiones relacionadas con tencologías 
y metodologías de seguridad en entornos académicos y de investigación. El \SMS\ proporcionó una base científica sólida para identificar 
tendencias, enfoques metodológicos, controles de seguridad y tecnologías ampliamente utilizadas en infraestructuras de computación en la 
nube, mientras que el análisis \DAR\ permitió evaluar de manera estructurada las alternativas tecnológicas mediante criterios técnicos, 
operativos y normativos claramente definidos.

Asimismo, la caracterización integral del grupo \GRID\ permitió comprender no solo las capacidades tecnológicas actuales de la 
infraestructura, sino también las necesidades operativas, problemáticas existentes y oportunidades de fortalecimiento en materia de 
seguridad. Este enfoque integral facilitó que las decisiones adoptadas se alinearan con los objetivos estratégicos del entorno 
institucional y con los requerimientos de protección, monitoreo y gestión de eventos definidos para la arquitectura propuesta.

\section{Sobre la selección e implementación de aspectos de seguridad}
\noindent

La aplicación del análisis \DAR\ permitió identificar de manera objetiva las tecnologías más adecuadas para fortalecer la arquitectura de 
seguridad del entorno \GRID, dando como resultado la selección de Suricata como solución \IDPS y Wazuh como plataforma \SIEM. La elección de 
estas herramientas respondió a criterios relacionados con capacidad operativa, interoperabilidad, documentación, escalabilidad, correlación 
de eventos y alineamiento con controles de las normas \ISO/IEC 27001 e \ISO/IEC 27017, permitiendo consolidar un ecosistema tecnológico 
coherente con las necesidades actuales de la infraestructura.

La integración conjunta de Suricata y Wazuh demostró la viabilidad de construir un modelo de defensa en profundidad basado en tecnologías 
interoperables y adaptables a entornos académicos con recursos limitados. Asimismo, los aspectos de seguridad definidos permitieron 
establecer capacidades de monitoreo continuo, detección de anomalías, centralización de registros y respuesta activa ante incidentes, 
fortaleciendo la protección de los activos de información del grupo \GRID\ sin requerir modificaciones disruptivas sobre la infraestructura 
tecnológica previamente existente.

\section{Sobre el impacto institucional y académico}
\noindent

La arquitectura de defensa en profundidad con la integración de las tecnologías y controles propuestos para el grupo \GRID\ trasciende el 
ámbito estrictamente técnico y genera impactos relevantes tanto a nivel institucional como académico. En el contexto de la docencia, la 
incorporación de mecanismos de monitoreo, correlación de eventos y supervisión continua fortalece las capacidades de protección de la 
infraestructura utilizada por estudiantes y docentes, proporcionando un entorno más seguro para el desarrollo de actividades académicas, 
investigativas y de experimentación tecnológica relacionadas con computación en la nube y ciberseguridad.

Desde la perspectiva investigativa, la arquitectura y sus complementos implementados permiten consolidar una base tecnológica orientada al 
análisis de eventos de seguridad, monitoreo de tráfico y gestión de incidentes, facilitando futuros trabajos relacionados con arquitecturas 
de seguridad, análisis de amenazas y evaluación de controles en entornos cloud. Asimismo, el modelado realizado mediante ArchiMate y la 
documentación estructurada del proyecto contribuyen a la reproducibilidad y transferencia de conocimiento, permitiendo que otros grupos 
académicos o institucionales adapten la arquitectura propuesta según sus propias necesidades operacionales y de seguridad.

\section{Sobre la escalabilidad y sostenibilidad}
\noindent

Los aspectos de seguridad propuestos fueron diseñados considerando la escalabilidad y adaptabilidad como principios fundamentales, 
permitiendo la incorporación futura de nuevos componentes, servicios o mecanismos de protección sin afectar la interoperabilidad general 
del entorno. La integración de herramientas como Suricata y Wazuh, junto con los mecanismos de segmentación, monitoreo y correlación de 
eventos implementados, validan la capacidad de la arquitectura para extender sus funciones de supervisión y respuesta conforme evolucionen 
las necesidades tecnológicas y de seguridad del grupo \GRID.

Asimismo, la sostenibilidad del proyecto se fortalece mediante la documentación técnica, el modelado arquitectónico realizado en ArchiMate 
y la estructuración reproducible de los componentes definidos dentro de la propuesta. Esto permite que otros entornos académicos o 
institucionales con necesidades similares puedan adaptar, implementar o escalar la arquitectura planteada, contribuyendo al fortalecimiento 
de prácticas de seguridad y monitoreo en infraestructuras de computación en la nube.

\section{Sobre las limitaciones identificadas}
\noindent

Durante el desarrollo del proyecto se identificaron limitaciones y oportunidades de mejora que pueden orientar trabajos futuros 
relacionados con los aspectos de seguridad implementados en la arquitectura. Aunque el modelo de defensa en profundidad implementado 
permite fortalecer las capacidades de monitoreo, correlación y respuesta ante incidentes, la infraestructura actual corresponde a un 
entorno prototipo adaptado a las capacidades tecnológicas disponibles dentro del grupo \GRID, por lo que futuras implementaciones podrían 
incorporar una mayor diversidad de dispositivos, mecanismos avanzados de automatización y entornos distribuidos de mayor escala.

Asimismo, aunque las herramientas seleccionadas cumplen adecuadamente con los objetivos planteados, existen posibilidades de ampliación 
relacionadas con capacidades avanzadas de análisis de comportamiento, automatización de respuestas, integración con servicios 
institucionales de autenticación y mecanismos predictivos basados en inteligencia artificial para la detección temprana de amenazas. Estas 
limitaciones no afectan la funcionalidad ni la validez de la arquitectura implementada, sino que representan oportunidades para la 
evolución continua de la infraestructura de seguridad propuesta.

\section{Reflexiones sobre la computación distribuida académica}
\noindent

El presente proyecto evidencia la importancia estratégica que tiene la implementación de aspectos de seguridad en entornos académicos 
y de investigación que operan sobre infraestructuras de computación en la nube. El fortalecimiento de capacidades de monitoreo, correlación 
de eventos y respuesta ante incidentes no solo contribuye directamente a la protección de los activos tecnológicos del grupo \GRID, sino 
que también favorece el desarrollo de entornos más confiables para actividades académicas, científicas y de experimentación tecnológica.

Asimismo, la experiencia desarrollada demuestra que tecnologías de seguridad como Suricata y Wazuh, ampliamente adoptados, especialmente 
en pequeñas y medianas empresas (PYMES) y en centros de operaciones de seguridad (SOC), pueden adaptarse satisfactoriamente a escenarios 
académicos con recursos limitados. Esto permite implementar mecanismos avanzados de supervisión y protección sin requerir inversiones 
excesivamente elevadas, facilitando la adopción de buenas prácticas alineadas con estándares internacionales de seguridad de la 
información como \ISO/IEC 27001 e \ISO/IEC 27017.

\section{Aportes al estado del arte}
\noindent

La presente investigación contribuye al fortalecimiento del estado del arte relacionado con aspectos de seguridad en entornos de 
computación en la nube aplicados al contexto académico. Desde una perspectiva metodológica, el trabajo proporciona un marco estructurado 
que integra caracterización contextual, estudio de mapeo sistemático de literatura, análisis \DAR\ y modelado arquitectónico, permitiendo 
establecer un proceso reproducible para la selección, evaluación e integración de tecnologías de seguridad en infraestructuras 
institucionales.

Desde el enfoque técnico, el proyecto documenta la integración y funcionamiento conjunto de herramientas como Suricata, Wazuh y SonicWall 
NSA 6600 dentro de una arquitectura de defensa en profundidad implementada sobre un entorno prototipo construido de acuerdo con las 
características y necesidades identificadas en la infraestructura del grupo \GRID. Asimismo, el modelado realizado mediante ArchiMate y la 
validación conceptual de los flujos operacionales permiten evidenciar la viabilidad de las soluciones propuestas, proporcionando 
configuraciones, estructuras y lineamientos que pueden ser adaptados o replicados en instituciones con necesidades similares de monitoreo, 
correlación de eventos y protección de infraestructuras cloud.

\section{Conclusión general}
\noindent

El proyecto desarrollado permitió cumplir satisfactoriamente los objetivos planteados, dando como resultado una arquitectura de defensa en 
profundidad de acuerdo a los aspectos de seguridad identificados, además de estar orientada al fortalecimiento de la supervisión, 
detección y respuesta ante incidentes de seguridad dentro del entorno \GRID. La integración de tecnologías como Suricata, Wazuh y 
SonicWall NSA 6600, junto con el modelado arquitectónico realizado mediante ArchiMate y Draw.io, permitió estructurar una propuesta 
alineada con controles de las normas \ISO/IEC 27001 e \ISO/IEC 27017, respaldada por documentación técnica y lineamientos de implementación 
reproducibles.

La metodología empleada, basada en la caracterización contextual del entorno, el estudio de mapeo sistemático de literatura, el análisis 
\DAR\ y el diseño arquitectónico, proporcionó un enfoque estructurado para la selección e integración de tecnologías de seguridad en 
infraestructuras académicas. Asimismo, los resultados obtenidos y las configuraciones propuestas constituyen una referencia aplicable a 
otros contextos institucionales con necesidades similares de monitoreo, correlación de eventos y protección de infraestructuras cloud.

Finalmente, la evidencia generada durante el desarrollo del proyecto demuestra que la implementación de arquitecturas de seguridad 
multicapa en entornos académicos contribuye significativamente al fortalecimiento de las capacidades de supervisión, gobernanza y 
protección de activos de información. Esto permite mejorar la visibilidad sobre los eventos de seguridad, fortalecer los mecanismos de 
respuesta ante incidentes y establecer bases técnicas para futuras ampliaciones de las capacidades de seguridad dentro del grupo \GRID\ y 
otros entornos de investigación similares.
